\documentclass[11pt,a4paper]{article}

% --------------------
% Packages
% --------------------
\usepackage{amsmath,amssymb}
\usepackage{graphicx}
\usepackage{subcaption}
\usepackage{booktabs}
\usepackage{float}
\usepackage{geometry}
\usepackage{hyperref}
\usepackage{enumitem}

\geometry{margin=1in}

% --------------------
% Title Information
% --------------------
\title{\textbf{CS-736: Mathematical Imaging and Image Reconstruction}\\
Assignment -- X-Ray Computed Tomography}
\author{
Arshit Singh\\
\texttt{(Roll Number / Entry Number)}}
\date{Instructor: Suyash P. Awate\\ \vspace{0.2cm} \today}

% --------------------
\begin{document}
\maketitle
\thispagestyle{empty}
\newpage

% ============================================================
\section*{Question 1: X-Ray Computed Tomography -- Radon Transform}
% ============================================================

\subsection*{Objective}

The objective of this problem is to implement the Radon transform of a 2D image by numerically approximating line integrals at multiple projection angles and offsets, and to analyze the effects of discretization and sampling choices on the resulting sinograms.

\vspace{0.3cm}

% ------------------------------------------------------------
\subsection*{(a) X-ray Line Integration}
% ------------------------------------------------------------

Given an image $f(x,y)$ defined on a discrete grid, the Radon transform is defined as
\[
\mathcal{R}f(t,\theta)
=
\int_{-\infty}^{\infty}
f(x,y)\,
\delta(x\cos\theta + y\sin\theta - t)\,dx\,dy.
\]

Using a parametric representation of the line
\[
\begin{aligned}
x(s) &= t\cos\theta - s\sin\theta, \\
y(s) &= t\sin\theta + s\cos\theta,
\end{aligned}
\]
the transform can be written as
\[
\mathcal{R}f(t,\theta)
=
\int_{-\infty}^{\infty}
f(x(s), y(s))\,ds.
\]

Since the image is discrete, this integral is approximated using a summation:
\[
\mathcal{R}f(t,\theta)
\approx
\sum_{k} f(x(s_k), y(s_k))\,\Delta s,
\]
where $\Delta s$ is the sampling step size along the ray.

\paragraph{Choice of $\Delta s$.}
The Radon transform is computed using $\Delta s = 0.5$, $1.0$, and $3.0$ pixel-width units. Smaller values provide better numerical accuracy at the expense of higher computational cost, while larger values lead to undersampling and approximation errors.

\paragraph{Interpolation scheme.}
Bilinear interpolation is used to evaluate image intensities at off-grid locations. This choice offers a good trade-off between accuracy and computational efficiency while preserving linearity of the transform. Pixels outside the image domain are assigned zero intensity.

% ------------------------------------------------------------
\subsection*{(b) Discrete Radon Transform}
% ------------------------------------------------------------

Using the integration function from part (a), the Radon transform is computed for the discrete sets
\[
t = -90, -85, \ldots, 85, 90,
\quad
\theta = 0, 5, 10, \ldots, 175.
\]

For each pair $(t,\theta)$, numerical integration is performed along the corresponding line to obtain a discrete sinogram.

\vspace{0.3cm}

% ---------- FIGURE PLACEHOLDER ----------
\begin{figure}[H]
\centering
\fbox{\parbox[c][6cm][c]{0.8\textwidth}{\centering
\textbf{Placeholder: Radon transform (sinogram)}\\
$\Delta s = 1.0$, $t \in [-90,90]$, $\theta \in [0,175]$
}}
\caption{Discrete Radon transform of the Shepp--Logan phantom.}
\label{fig:radon_basic}
\end{figure}
% --------------------------------------

% ------------------------------------------------------------
\subsection*{(c) Effect of Integration Step Size}
% ------------------------------------------------------------

Radon-transform images are computed for three different step sizes: $\Delta s = 0.5$, $1.0$, and $3.0$ pixel-width units.

\paragraph{1D Radon profiles.}
For $\theta = 0^\circ$ and $\theta = 90^\circ$, the 1D Radon profiles computed with $\Delta s = 0.5$ appear the smoothest, while those computed with $\Delta s = 3.0$ appear the roughest due to undersampling along the integration direction.

% ---------- FIGURE PLACEHOLDER ----------
\begin{figure}[H]
\centering
\fbox{\parbox[c][6cm][c]{0.85\textwidth}{\centering
\textbf{Placeholder: 1D Radon profiles}\\
$\theta = 0^\circ$ and $90^\circ$ for different $\Delta s$
}}
\caption{Comparison of 1D Radon-transform profiles for different integration step sizes.}
\label{fig:radon_1d}
\end{figure}
% --------------------------------------

\paragraph{2D Radon-transform images.}
The Radon-transform image corresponding to $\Delta s = 0.5$ is the smoothest, while the image corresponding to $\Delta s = 3.0$ shows visible artifacts and loss of detail.

% ---------- FIGURE PLACEHOLDER ----------
\begin{figure}[H]
\centering
\fbox{\parbox[c][6cm][c]{0.9\textwidth}{\centering
\textbf{Placeholder: Radon images for}\\
$\Delta s = 0.5,\; 1.0,\; 3.0$
}}
\caption{Effect of integration step size on the Radon-transform image.}
\label{fig:radon_ds}
\end{figure}
% --------------------------------------

% ------------------------------------------------------------
\subsection*{(d) CT Scanner Parameter Design}
% ------------------------------------------------------------

In designing an X-ray CT scanner, smaller values of $\Delta t$ and $\Delta \theta$ are preferred to ensure dense spatial and angular sampling. Smaller $\Delta t$ improves spatial resolution, while smaller $\Delta \theta$ reduces angular undersampling artifacts. However, smaller values increase acquisition time, radiation dose, and storage requirements, leading to a trade-off between image quality and practical constraints.

% ------------------------------------------------------------
\subsection*{(e) ART Reconstruction Considerations}
% ------------------------------------------------------------

\paragraph{Choice of pixel grid.}
The pixel size should be chosen small enough to capture anatomical details but large enough to avoid an excessively ill-conditioned reconstruction problem.

\paragraph{Effect of $\Delta s$.}
Choosing $\Delta s \gg 1$ pixel width leads to inaccurate ray–pixel interactions and degraded reconstructions, while choosing $\Delta s \ll 1$ pixel width improves accuracy but increases computational cost with diminishing returns.

% ============================================================
\end{document}
